% Section declaración - SECTOR FINANCIERO
\par{

	Economista con sólida formación en finanzas, análisis cuantitativo y gestión de riesgos. Apasionado por el sector financiero y su rol en el desarrollo económico. Con capacidad analítica demostrada, dominio de herramientas estadísticas y conocimiento profundo del sistema financiero peruano e internacional.

	Experiencia en análisis de indicadores financieros, evaluación de riesgo crediticio, procesamiento de datos financieros y elaboración de reportes regulatorios. Familiarizado con normativa SBS, estándares de Basilea, análisis de estados financieros y productos financieros (créditos, inversiones, seguros).

	Busco la oportunidad de integrarme a una institución financiera sólida y en crecimiento donde pueda aplicar mis conocimientos en economía y finanzas, desarrollar mi carrera profesional en un entorno dinámico y desafiante, y contribuir al cumplimiento de objetivos institucionales con profesionalismo, ética y orientación a resultados.

	%%%%%%%%%%%%%%%%%%%%%%%%%%%%%%%%%%%%%%%%
	% VERSIÓN ALTERNATIVA - Risk Management
	%%%%%%%%%%%%%%%%%%%%%%%%%%%%%%%%%%%%%%%%
	%Economista con especialización en gestión de riesgos financieros y análisis cuantitativo. Experto en modelamiento de riesgo crediticio, construcción de scorecards, análisis de morosidad y proyecciones de pérdidas esperadas.

	%Dominio de metodologías de credit scoring, análisis de vintage, matrices de transición y backtesting de modelos. Capacidad de trabajar con bases de datos de clientes, variables macroeconómicas y desarrollar modelos predictivos con R, Python y software estadístico especializado.

	%Busco oportunidades en áreas de riesgos, credit risk modeling o analytics financieros donde pueda aplicar mi formación cuantitativa para desarrollar modelos robustos, cumplir con requerimientos regulatorios y contribuir a la toma de decisiones de crédito basada en evidencia y mejores prácticas internacionales.

	%%%%%%%%%%%%%%%%%%%%%%%%%%%%%%%%%%%%%%%%
	% VERSIÓN ALTERNATIVA - Análisis financiero
	%%%%%%%%%%%%%%%%%%%%%%%%%%%%%%%%%%%%%%%%
	%Economista con orientación en análisis financiero y evaluación de inversiones. Especializado en análisis de estados financieros, valorización de empresas, evaluación de proyectos de inversión y análisis de mercados de capitales.

	%Dominio de herramientas de valorización (DCF, múltiplos), análisis de ratios financieros, construcción de modelos en Excel y análisis de sensibilidad. Familiarizado con Bloomberg, Reuters y fuentes de información financiera. Capacidad de elaborar informes de análisis fundamental y recomendaciones de inversión.

	%Busco oportunidades en banca de inversión, asset management, equity research o áreas de finanzas corporativas donde pueda aplicar mi capacidad analítica y mi conocimiento de mercados financieros para evaluar oportunidades, generar alpha y contribuir a la creación de valor para inversionistas y accionistas.

	%%%%%%%%%%%%%%%%%%%%%%%%%%%%%%%%%%%%%%%%
	% VERSIÓN ALTERNATIVA - Productos financieros
	%%%%%%%%%%%%%%%%%%%%%%%%%%%%%%%%%%%%%%%%
	%Economista con conocimiento profundo del ecosistema de productos financieros y foco en experiencia del cliente. Especializado en análisis de demanda, diseño de productos, estrategias de pricing y segmentación de mercados financieros.

	%Experiencia en análisis de comportamiento de clientes, estudios de mercado financiero, análisis de competencia y propuestas de nuevos productos o mejoras. Capacidad de trabajar con equipos comerciales, tecnología y compliance para desarrollar soluciones financieras competitivas y reguladas.

	%Busco oportunidades en product management, desarrollo de negocios o inteligencia comercial en el sector financiero donde pueda combinar mi visión analítica con orientación al cliente para impulsar la innovación, rentabilidad y crecimiento sostenible del portafolio de productos financieros.

	%%%%%%%%%%%%%%%%%%%%%%%%%%%%%%%%%%%%%%%%
	% VERSIÓN ALTERNATIVA - Microfinanzas y banca inclusiva
	%%%%%%%%%%%%%%%%%%%%%%%%%%%%%%%%%%%%%%%%
	%Economista con vocación social y especialización en microfinanzas, inclusión financiera y finanzas populares. Experto en análisis de mercados de base de la pirámide, evaluación de tecnologías crediticias alternativas y diseño de productos financieros para segmentos no bancarizados.

	%Experiencia en análisis de pobreza, informalidad, características socioeconómicas de micro y pequeños empresarios. Familiarizado con metodologías de crédito grupal, banca comunal, credit scoring alternativo y canales digitales para inclusión financiera.

	%Busco oportunidades en instituciones microfinancieras, ONGs de desarrollo o áreas de inclusión financiera donde pueda aplicar mi formación económica y mi compromiso social para expandir el acceso a servicios financieros de calidad, promover el desarrollo económico de poblaciones vulnerables y contribuir a la reducción de la pobreza.

	%%%%%%%%%%%%%%%%%%%%%%%%%%%%%%%%%%%%%%%%
	% VERSIÓN ALTERNATIVA - Finanzas corporativas
	%%%%%%%%%%%%%%%%%%%%%%%%%%%%%%%%%%%%%%%%
	%Economista con especialización en finanzas corporativas y planeamiento financiero. Experto en elaboración de presupuestos, proyecciones financieras, análisis de estructura de capital y evaluación de decisiones de financiamiento e inversión.

	%Dominio de modelamiento financiero en Excel, análisis de flujos de caja, cálculo de WACC, evaluación de VAN/TIR y análisis de punto de equilibrio. Capacidad de elaborar business plans, pitch decks para inversionistas y reportes financieros ejecutivos.

	%Busco oportunidades en áreas de finanzas corporativas, FP&A (Financial Planning & Analysis) o tesorería donde pueda aplicar mi visión estratégica y mis habilidades cuantitativas para optimizar la estructura financiera, maximizar el valor de la empresa y apoyar decisiones de inversión y financiamiento con análisis rigurosos.

}
